\documentclass[10pt,a4paper]{amsart}
\usepackage[margin=19mm]{geometry}
\usepackage{amsmath,amssymb,mathtools}
\usepackage[scr=rsfs,cal=euler]{mathalfa}
\usepackage{xfrac}
\usepackage[unicode,hidelinks,bookmarks=false]{hyperref}
\newcommand{\ie}{i.e.\ }
\newcommand{\cf}{cf.\ }
\newcommand{\resp}{respectively\ }
\newcommand{\Subsection}{\subsection}
\providecommand{\nobreakdash}{\nobreak\hbox{-}}
\setlength{\emergencystretch}{2em}
\multlinegap0pt
\usepackage{bm}
\usepackage{bbm}
\usepackage{dsfont}
\usepackage{xspace}
\usepackage{cancel}
\usepackage{mathtools}
\usepackage{amsthm}
\makeatletter
\@ifundefined{theorem}{\newtheorem{theorem}{Theorem}}{}
\@ifundefined{corollary}{\newtheorem{corollary}[theorem]{Corollary}}{}
\@ifundefined{definition}{\newtheorem{definition}[theorem]{Definition}}{}
\@ifundefined{lemma}{\newtheorem{lemma}[theorem]{Lemma}}{}
\@ifundefined{proposition}{\newtheorem{proposition}[theorem]{Proposition}}{}
\makeatother
\def\B{\mathcal{B}}
\def\G{\mathcal{G}}
\def\L{\mathcal{L}}
\def\P{\mathcal{P}}
\def\W{\mathcal{W}}
\def\Y{\mathcal{Y}}
\def\ref#1{#1}
\newcommand{\X}{mathcal{X}}
\title[Conjectural Positivity for Pontryagin Product in Equivariant]{Conjectural Positivity for Pontryagin Product in Equivariant K-theory of Loop Groups}
\author{Shrawan Kumar}
\date{Source: arXiv:2510.07689v1 (CC BY 4.0)}
\begin{document}
\maketitle

\noindent\emph{Source and licence.} Conjectural Positivity for Pontryagin Product in Equivariant K-theory of Loop Groups. Version arXiv:2510.07689v1 (CC BY 4.0), distributed under the Creative Commons Attribution 4.0 International licence, as recorded in the arXiv licence metadata for this exact version and shown on the abstract page https://arxiv.org/abs/2510.07689v1. Full rights notice: licenses/arxiv-2510.07689-CC-BY-4.0.txt.\par\smallskip

\noindent\textit{Editorial scope.} Counted unit: the Proposition whose source label is \texttt{6.3} in the article (source0.tex lines 1860--1868, source label \texttt{6.3}), delivered together with the article's own complete original proof of it (source0.tex lines 1870--1918, one \texttt{proof} environment of 49 lines). No mathematics is written, repaired or re-derived: statement and proof are the article's own text line for line. Required context retained uncounted: 4.2 (lines 809--822); 4.3 (lines 824--832); 4.5 (lines 940--950); 4.7 (lines 1002--1007); 4.8 (lines 1033--1110); 4.8.4 (lines 1090--1092); 4.11 (lines 1139--1141); 5.9 (lines 1343--1350); 6.1 (lines 1832--1843). Every definition, display and label of those lines is retained unchanged. Omitted from this package and not redistributed: 1--808, 823--823, 833--939, 951--1001, 1008--1032, 1093--1138, 1142--1342, 1351--1831, 1844--1859, 1869--1869, 1919--2245. Nothing inside a retained excerpt is omitted or altered: every retained body is delivered exactly as the article's lines are written, so no omission receipt of any kind accompanies this package. Citations: the retained excerpts carry no citation command at all, so no bibliography entry of the article is transcribed and no reference is redistributed. Reference adaptation: none was needed and none was made; every reference inside the delivered unit resolves inside it, so no unresolved reference or citation is produced. Printed slips kept literal: no printed word, symbol, hypothesis or number of the counted statement or of its proof was altered. Layout and class: the article's own class is \texttt{article} with packages including amsthm, bbm, amsmath, MnSymbol, todonotes, mathtools, mathrsfs, xy, nicefrac, esvect, babel, color, hyperref, fontenc; the wrapper is the lane's pinned \texttt{amsart} editorial wrapper at 10pt on A4, which restates the theorem-like environments the retained text needs and adds the article's own macro definitions that the retained text needs (\texttt{\textbackslash B}, \texttt{\textbackslash G}, \texttt{\textbackslash L}, \texttt{\textbackslash P}, \texttt{\textbackslash W}, \texttt{\textbackslash X}, \texttt{\textbackslash Y}, \texttt{\textbackslash ref}), with extra wrapper packages (bm, bbm, dsfont, xspace, cancel, mathtools). The delivered numbering is the unit's own. Assets: no retained span contains a figure environment, a graphics-inclusion command, a PDF-page inclusion, a table or a TikZ picture; no image file is copied into this package, the package declares no asset dependency and no image file is redistributed with it. Licence: version arXiv:2510.07689v1 of this article is distributed under the Creative Commons Attribution 4.0 International licence, as recorded in the arXiv licence metadata for this exact version and shown on the abstract page https://arxiv.org/abs/2510.07689v1. A full rights notice is delivered at licenses/arxiv-2510.07689-CC-BY-4.0.txt. Characters: every retained excerpt is pure ASCII, so the delivered engine, XeLaTeX with its default Latin Modern text font, reports no missing character.
\begin{definition}\label{4.2}
\rm{
In any Coxeter group $\W$,  define the {\it Demazure product} $*$ for any $u\in \W$ and simple reflection $s_i$,
$$
u* s_i=
\left\{ \begin{array}{l}
u,\,\,\,  \mbox{if}\quad us_i <u\\
us_i,\,\,\, \mbox{if}\quad us_i >u.
\end{array}\right.
$$
This extends to an associative product by defining 
$$u*v= (\cdots (( u*s_{i_1})*s_{i_2})\cdots *s_{i_n})$$
 for a reduced decomposition $v=s_{i_1}\ldots s_{i_n}$. (It does not depend upon the choice of the reduced decomposition of $v$.)}
\end{definition}

\begin{proposition} \label{4.3}
For $u\in \W$ and $v\in \W'$,
$$
\left[\mathcal{O}_{X^\B_u}\right] \odot'\left [\mathcal{O}_{X_v}\right] =\left [\mathcal{O}_{X_{\overline{u*v}}}\right] \in K_0^B\left(\X\right) ,
$$
where ${X^\B_u}:= \overline{\B u\B/\B} \subset \Y$.

Observe that $u*v$ may not lie in $\W'$. We take its unique representative $\overline{u*v}$ in $\W'$. 
\end{proposition}

\begin{corollary} \label{4.5}
For  $b_0 \in K_0^\P (*) = K_0^G (*)$,
$$
\left(\mu_i\right)_! \left( \mathcal{O}_{X_i} \boxtimes^\B b_0\right) =b_0 \left[\mathcal{O}_{X_i}\right] \in K_0^\B\left(X_i\right).
$$

Thus, following the proof of Proposition \ref{4.3}, we get that for any $a\in K_0^\B\left(\Y \right)$ and $b\in K_0^\B\left(\X\right)$,
$$
a\odot' \left(b_0\cdot b\right) = b_0a\odot' b. \qquad\qquad\qquad \square
$$
\end{corollary} 

\begin{lemma} \label{4.7}
With the notation as above
$$
\Delta =\sum_{x\in W}\L_x \boxtimes \L^x \in K_0^\P \left(\P/\B \times \P/\B\right). 
$$
\end{lemma}

\begin{definition}\label{4.8} \rm{
Consider the commutative diagram:
$$
\arraycolsep=2pt
\begin{array}{lcl}
\displaystyle K_0^\B\left(\X\right)&\displaystyle \overset{\overset{i}{\sim}}{\longrightarrow} K_0^\P \left(\P \times^\B \X\right) &\displaystyle \underset{\sim}{\overset{\eta}{\longrightarrow}} K_0^{\P} \left(\P/\B \times \X\right)\\
\rotatebox[origin=c]{140}{$\underset{\sim}{\overset{\underline{\phi}}{\longleftarrow}}$} 
&& \rotatebox[origin=c]{240}{$\underset{\sim}{\overset{{\phi}}{\longrightarrow}}$}  
\\
&\displaystyle  K_0^\P \left(\P/\B\right)  \underset{K_0^\P(*)}{\boxtimes} K_0^\P\left(\X\right).&
\end{array}
$$
In this diagram $i$ is the Induction Isomorphism [CG, \S5.2.15], the isomorphism $\eta$ is induced from the $\P$-equivariant isomorphism of the ind-varieties:
$$
\P \times^{\B} \X \underset{\sim}{\rightarrow} \P/\B \times  \X, \quad \left[ p,x\right] \mapsto\left(p\B,px\right),
\qquad \mbox{for } p\in \P \mbox{ and } x\in \X.
$$

The isomorphism $\phi$ is the Kunneth isomorphism (cf. [CG, Theorem 5.6.1]). To satisfy the hypotheses of loc cit., we have used Lemma \ref{4.7} and the result that 
$$ 
K_0^\P \left(\Y\right) =K_0^G \left(\Y\right),\quad\mbox{ for any } \mbox{ $\P$-ind-variety }\Y. 
$$

By definition, $\overline{\phi}=\phi\circ \eta \circ i$ and hence it is an isomorphism. Analyzing the proof of [CG, Theorem 5.6.1], specifically on page 275 of loc cit., we get that 
\begin{equation}\label{4.8.1}
\overline{\phi} (b)=\sum_{x\in W} \L^x \boxtimes \overline{\mu}_! \left( \L_x \overset{\B}{\boxtimes} b\right), \mbox{ for any } b\in K_0^\B \left(\X\right),
\end{equation}
where $\overline{\mu}:\P\times^\B \X\rightarrow \X$ is the product map $[p,x]\mapsto p\cdot x$, for $p\in\P$ and $x\in \X$. Here, we have  abbreviated $\left(\hat{p}^* \L_x \right) \overset{\B}{\boxtimes} b$
by  $\L_x\overset{\B}{\boxtimes} b$, where $\hat{p}:\P\rightarrow \P/\B$ is the projection.
In particular,  for $b=\mathcal{O}_{X_u}$ (for $u\in \W'$) ,
\begin{equation}\label{4.8.2}
\overline{\phi} \left(\left[ \mathcal{O}_{X_u}\right]\right) = \sum_{x\in W} \L^x \boxtimes \overline{\mu}_! \left( \L_x \overset{\B}{\boxtimes}  \mathcal{O}_{X_u}\right). 
\end{equation}
As mentioned earlier, $\odot'$ is \textit{not} $R(\B)$-linear in the second variable. To remedy this, we modify its definition following [Ka-2, \S8].
Define the \textit{modified convolution product}: 
$$
\odot :K_0^\B \left( \Y\right) \bigotimes _{K_0^\B(*)} K_0^\B \left( \X\right) \rightarrow K_0^\B \left ( \X\right)
$$
by 
\begin{equation}\label{4.8.3}
a\odot b := \sum_{x\in W} \left(\epsilon \left(\L^x\right) \cdot a\right) \, \odot' \overline{\mu}_! \left(\L_x \overset{\B}{\boxtimes} b\right), \mbox{ for  } a\in K_0^\B\left(\Y\right) \mbox{ and } b \in K_0^\B\left(\X\right),
\end{equation}
where $\epsilon:K_0^\P \left(\P/\B\right)\overset{\sim}{\rightarrow}K_0^\B (*)$ is the isomorphism $i^{-1}$ as earlier  for $\X$ replaced by $*$. It is easy to see that $\odot$ does not depend on the choice of the basis $\L_w$. From the definition of $\odot$, it follows that $\odot$ is $K_0^\B(*)$-bilinear. It is clearly $K_0^\B(*)$-linear in the first variable. To prove its linearity in the second variable, take a character $e^\lambda$ of $\B$. Then, 
\begin{eqnarray*}
a\odot e^\lambda \cdot b &=& \sum_{x\in W}\epsilon \left(\L^x\right)\cdot a \odot' \overline{\mu}_!\left(\L_x \overset{\B}{\boxtimes} e^\lambda \cdot b\right)\\
&=& \sum_{x\in W}\epsilon \left(\L^x\right)\cdot a \odot' \overline{\mu}_!\left(\L(-\lambda)\cdot \L_x \overset{\B}{\boxtimes} b\right)\\
&=& \sum_{x,y\in W}\epsilon \left(\L^x\right) \left\langle \L(-\lambda)\cdot \L_x,\L^y\right\rangle\cdot a \odot' \overline{\mu}_! \left( \L_y \overset{\B}{\boxtimes} b\right)\\
&&\text{since $\odot'$ is $R(\P)$-linear in the second variable by Corollary \ref{4.5}}\\
&=&\sum_{y\in W}\left( \sum_{x\in W} \epsilon \left(\L^x\right) \left\langle \L_x, \L\left(-\lambda\right)\L^y \right\rangle\right) \cdot a \odot' \overline{\mu}_! \left( \L_y \overset{\B}{\boxtimes} b\right)\\
&=&\sum_{y\in W} \epsilon\left(\L\left(-\lambda\right)\L^y\right)\cdot a \odot' \overline{\mu}_! \left( \L_y \overset{\B}{\boxtimes} b\right) ,\\
&&\quad\mbox{since }\sum_{x\in W}\L^x \left\langle \L_x, \L\left(-\lambda\right)\L^y \right\rangle = \L\left(-\lambda\right)\L^y\\
&=&\sum_{y\in W}e^\lambda \cdot \epsilon \left(\L^y\right) \cdot a \odot'  \overline{\mu}_! \left( \L_y \overset{\B}{\boxtimes} b\right)\\
&=&e^\lambda a \odot b.
\end{eqnarray*}
This proves that $\odot$ is $K_0^\B(*)$-linear in the second variable. 

We now prove that for $b\in K_0^\P(\X)$, 
\begin{equation}\label{4.8.4}
a\odot b= a\odot' b, \quad \mbox{for any } a\in K_0^\B \left(\Y\right).
\end{equation}
Since $b\in K_0^\P(\X)$, it is easy to see that, for any $x\in W$, 
$$
 \overline{\mu}_! \left( \L_x\overset{\B}{\boxtimes} b\right)= \chi_\P\left(\L_x\right)\cdot b,\,\,\,\text{by the projection formula,}
 $$ since $\L_x\overset{\B}{\boxtimes} b=\tilde{\pi}^* \left(\L_x\right)\otimes \left(\epsilon \overset{\B}{\boxtimes} b\right)$, where $\tilde{\pi}: \P\times^\B \X \rightarrow \P/\B$ is the projection. 
 
 Thus, 
\begin{eqnarray*}
a\odot b &=&  \sum_{x\in W} \epsilon \left(\L^x\right) \cdot a \odot' \chi_\P\left(\L_x\right)\cdot b\\
&=& \sum_{x\in W} \epsilon \left(\L^x\right) \chi_\P\left(\L_x\right)\cdot a \odot' b,\\
&& \mbox{since $\odot'$ is $K_0^\P(*)=R(\P)$-linear in the second variable }\\
&=&\epsilon \left(\mathcal{O}_{\P/\B}\right)\cdot a \odot' b,\mbox{ as above since } \,\chi_\P(\L_x) :=\langle \L_x, \mathcal{O}_{\P/\B}\rangle\\
&=&a \odot' b.
\end{eqnarray*}
This proves \eqref{4.8.4}. 


Let $*$ be the Pontryagin product in $K_0^T(\X)$ as in Definition 3.1 and $\odot$ the modified convolution product $K_0^T(\Y)\otimes K_0^T(\X)\rightarrow K_0^T(\X)$ as above. Since $p:\Y \rightarrow \X$ is a $\G$-equivariant (in particular, $\B$-equivariant) fibration; in particular, it is a flat morphism. Thus, there is the pull-back map $p^*:K_0^T(\X) \rightarrow K_0^T(\Y)$. This takes, for $w\in \W'$, $[\mathcal{O}_{X_w}]\mapsto[\mathcal{O}_{X^\B_{xw_o}}]$, where $w_o$ is the longest element of $W$. Via this $p^*$, we get a (modified) convolution product $\odot$ on $K_0^T(\X)$. }
\end{definition}

\begin{corollary}\label{4.11}
The product $\odot$ in $K_0^T(\X)$ is associative and commutative. 
\end{corollary}

\begin{theorem} \label{5.9}
Take $u\in \W$, $v\in \W'$ and take a reduced decomposition $u=s_{i_1}\ldots s_{i_n}$\,  $\left(0\leq i_j\leq l\right)$. Then, under the modified convolution product (cf. Definition \ref{4.8})
$$
\left[ \mathcal{O}_{X^\B_u}\right] \odot \left[ \mathcal{O}_{X_v}\right] = \sum_{x\in W} \sum_{1\leq j_1< \cdots <j_p\leq n} \left\vert D_{i_1}' \cdots \hat{\hat{D}}'_{i_{j_1}}\cdots \hat{\hat{D}}'_{i_{j_p}} \cdots D_{i_n}'\left(\bar{\zeta}^x \right)\right\vert [ \mathcal{O}_{X_{\overline{s_{i_{j_1}}*\cdots *s_{i_{j_p}}*x*v}}}],
$$
where $\hat{\hat{D}}'_j$ means to replace the operator $D_j'$ by the Weyl
group action on $R(T) \otimes_{R(G)} R(T)$ acting only on the first factor, $*$ is the Demazure product in $\W$ (cf. Definition \ref{4.2}) and for $w\in \W$, $\bar{w}$ denotes the corresponding minimal representative in $wW$. 
\end{theorem}

\begin{lemma}\label{6.1}
(a) $\bar\zeta^e=e^{-\rho}\otimes e^\rho$, $\bar\zeta^{s_1}=e^0\otimes e^0-e^{-\rho}\otimes e^\rho$ ,  where $\rho=\frac{\alpha_1}{2}$ .
\vskip1ex

(b) $s_0'\left(\bar\zeta^e\right)=e^\rho\otimes e^\rho$, $s_0'\left( \bar\zeta^{s_1}\right)=\bar\zeta^{s_1}+\left(1-e^{2\rho}\right)\bar\zeta^e$ .
\vskip1ex

(c) $D'_0\left(\bar\zeta^e\right)=-e^{2\rho}\bar\zeta^e$, $D'_0\left(\bar\zeta^{s_1}\right)=e^{2\rho}\bar\zeta^e$ .
\vskip1ex

(d) $D'_1\left(\bar\zeta^e\right)=\bar\zeta^e$, $D'_1\left(\bar\zeta^{s_1}\right)=-\bar\zeta^e$ .
\end{lemma}

\begin{proposition} \label{6.3}
For any $n,m\geq0$ , under the (modified) convolution product $\odot$ in $K_0^T(\X)$, 
\vskip1ex
(a) $\left[ \mathcal{O}_n\right] \odot \left[ \mathcal{O}_{2m}\right]=\left[ \mathcal{O}_{n+2m}\right]  $
\vskip1ex
(b) $\left[ \mathcal{O}_{2n+1}\right] \odot \left[ \mathcal{O}_{2m+1}\right]=e^{\alpha_1}\left[ \mathcal{O}_{2n+2m+2}\right]+\left(1-e^{\alpha_1}\right) \left[ \mathcal{O}_{2n+2m+3}\right]$,

where $\mathcal{O}_n$ denotes $\mathcal{O}_{X_n}$. 
\end{proposition} 

\begin{proof}
We first calculate for $v=\tau_{2m}$ (denoting $\mathcal{O}_v=\mathcal{O}_{X_v}$)
\begin{eqnarray*}
\left[ \mathcal{O}_{1}\right] \odot \left[ \mathcal{O}_{2m}\right]&=& \left[ \mathcal{O}_{X_{s_0s_1}^\B}\right] \odot \left[ \mathcal{O}_{\tau_{2m}}\right]\\
&=& \left\vert s_0's_1'\left( \bar{\zeta}^e+\bar{\zeta}^{s_1}\right) \right\vert \left[\mathcal{O}_{s_0s_1*\tau_{2m}}\right]
+ \left\vert D_0's_1' \left(\bar{\zeta}^e+\bar{\zeta}^{s_1}\right) \right\vert \left[\mathcal{O}_{s_1*\tau_{2m}}\right]\\
&&+ \left\vert s_0'D_1' \left(\bar{\zeta}^e+\bar{\zeta}^{s_1}\right) \right\vert \left[\mathcal{O}_{s_0*\tau_{2m}}\right]\\
&&+ \left\vert D_0'D_1' \left(\bar{\zeta}^e+\bar{\zeta}^{s_1}\right) \right\vert \left[\mathcal{O}_{\tau_{2m}}\right],\mbox{ by Theorem \ref{5.9}}\\
&=&  \left[\mathcal{O}_{s_0*\tau_{2m}}\right],\mbox{ by Lemma \ref{6.1}(a), since } \bar{\zeta}^e+\bar{\zeta}^{s_1}= e^0\otimes e^0. 
\end{eqnarray*}
Thus, 
\begin{equation}\label{6.3.1}
\left[ \mathcal{O}_{1}\right] \odot \left[ \mathcal{O}_{2m}\right] = \left[ \mathcal{O}_{2m+1}\right].
\end{equation}
We next calculate  
\begin{eqnarray*}
\left[ \mathcal{O}_{1}\right] \odot \left[ \mathcal{O}_{2m+1}\right] &=& \left[ \mathcal{O}_{X_{s_0s_1}^\B}\right] \odot \left[ \mathcal{O}_{2m+1}\right] \\
&=&  \left\vert D_0'D_1' \left(\bar{\zeta}^e\right) \right\vert \left[\mathcal{O}_{2m+1}\right]+  \left\vert D_0'D_1' \left(\bar{\zeta}^{s_1}\right) \right\vert \left[\mathcal{O}_{s_1*\tau_{2m+1}}\right]\\
&&+ \left\vert D_0's_1' \left(\bar{\zeta}^e+\bar{\zeta}^{s_1}\right) \right\vert \left[\mathcal{O}_{s_1*\tau_{2m+1}}\right]+\left\vert s_0'D_1'(\bar{\zeta}^e)\right\vert \left[\mathcal{O}_{s_0*\tau_{2m+1}}\right]\\
&&+ \left\vert s_0'D_1' (\bar{\zeta}^{s_1}) \right\vert \left[\mathcal{O}_{s_0*s_1*\tau_{2m+1}}\right]\\
&&+\left\vert s_0's_1'\left(\bar{\zeta}^e+\bar{\zeta}^{s_1}\right) \right\vert \left[\mathcal{O}_{s_0*s_1*\tau_{2m+1}}\right],\mbox{ by Theorem \ref{5.9}}\\
&=& -e^{2\rho} \left[\mathcal{O}_{2m+1}\right] +e^{2\rho} \left[\mathcal{O}_{2m+2}\right] + e^{2\rho} \left[\mathcal{O}_{2m+1}\right] \\
&&- e^{2\rho} \left[\mathcal{O}_{2m+3}\right] +\left[\mathcal{O}_{2m+3}\right] \mbox{, by Lemma \ref{6.1} }.
\end{eqnarray*}
Thus, 
\begin{equation}\label{6.3.2}
\left[ \mathcal{O}_{1}\right] \odot \left[ \mathcal{O}_{2m+1}\right] = e^{2\rho} \left[\mathcal{O}_{2m+2}\right] +\left(1-e^{2\rho} \right)\left[\mathcal{O}_{2m+3}\right].
\end{equation}
Similar to the calculation of equation \eqref{6.3.1}, for any $n\geq 0$, writing $\tau_n\cdot s_1=s_{i_1}\ldots s_{i_{n+1}}$ with $s_{i_j}\in \left\{ 0,1\right\}$ and $s_{i_{n+1}}=s_1$ , we get
\begin{eqnarray*}
\left[ \mathcal{O}_{n}\right] \odot \left[ \mathcal{O}_{2m}\right]  &=& \left[ \mathcal{O}_{X_{\tau_n,s_1}^\B}\right] \odot \left[\mathcal{O}_{\tau_{2m}}\right]\\
&=& \sum_{1\leq j_1 <\cdots < j_p \leq n+1}\left\vert D_{i_1}' \circ \cdots \circ \hat{\hat{D}}'_{i_{j_1}}\circ \cdots \circ \hat{\hat{D}}'_{i_{j_p}}\circ \cdots \circ D_{i_{n+1}}'\left(\bar{\zeta}^e+\bar{\zeta}^{s_1}\right) \right\vert \\
&&[ \mathcal{O}_{X_{\overline{s_{i_{j_1}}*\cdots * s_{i_{j_p}}*\tau_{2m}}}}],\mbox{ by Theorem \ref{5.9}} .
\end{eqnarray*}
Thus,
\begin{equation}\label{6.3.3}
\left[ \mathcal{O}_{n}\right] \odot \left[ \mathcal{O}_{2m}\right] =  \left[ \mathcal{O}_{n+2m}\right] ,\mbox{ since } \bar{\zeta}^e+\bar{\zeta}^{s_1} = e^0 \otimes e^0.
\end{equation}
From equation \eqref{6.3.3}, we get (a). 

To prove (b), from (a) we get,
\begin{eqnarray*}
\left[ \mathcal{O}_{2n+1}\right] \odot \left[ \mathcal{O}_{2m+1}\right]  &=& \left(\left[ \mathcal{O}_{1}\right] \odot \left[ \mathcal{O}_{2n}\right] \right)\odot  \left[ \mathcal{O}_{2m+1}\right] \\
&=& \left[ \mathcal{O}_{1}\right] \odot \left[ \mathcal{O}_{2n+2m+1}\right] \mbox{, from the associativity}\\
&&\mbox{ and commutativity of $\odot$ as in  Corollary \ref{4.11}}\\
&=&e^{2\rho}\left[ \mathcal{O}_{2n+2m+2}\right] +\left(1-e^{2\rho}\right)\left[ \mathcal{O}_{2n+2m+3}\right], \mbox{  by \eqref{6.3.2}.}
\end{eqnarray*}
This proves (b).
\end{proof}

\end{document}
